\section{Objectives}

The primary objective of this project is to design and develop an integrated system that combines the ESP32 TTGO T-Display microcontroller with a custom 3D-printed housing. This project aims to explore the potential of additive manufacturing in creating functional and aesthetically pleasing enclosures for electronic devices, while also applying design for assembly principles to ensure ease of manufacturing and maintenance. The specific objectives of the project include:

\begin{itemize}
    \item To design a custom housing for the ESP32 TTGO T-Display using 3D modeling software, ensuring that it accommodates the device's components and allows for easy access to ports and buttons.
    \item To optimize the design of the housing for 3D printing, considering factors such as material selection, print orientation, and support structures to achieve a high-quality finish and structural integrity.
    \item To apply design for assembly principles in the housing design to facilitate easy assembly and disassembly of the device, allowing for maintenance and potential upgrades.
    \item To prototype the designed housing using a 3D printer and test its fit and functionality with the ESP32 TTGO T-Display, making necessary adjustments based on testing results.
    \item To document the design process, including challenges faced and solutions implemented, to provide insights into the integration of 3D printing with electronic device design.
    \item To demonstrate the potential of combining 3D printing with electronic device design to create innovative and functional products that can be easily manufactured and maintained.
\end{itemize}

Through the successful completion of these objectives, this project aims to contribute to the understanding of how additive manufacturing can be effectively utilized in the design and production of electronic device housings, while also showcasing the benefits of design for assembly in enhancing the manufacturability and maintainability of such devices.